\section{Metadatos}\label{metadatos}

{\def\LTcaptype{none} % do not increment counter
\begin{longtable}[]{@{}ll@{}}
\toprule\noalign{}
Campo & Valor \\
\midrule\noalign{}
\endhead
\bottomrule\noalign{}
\endlastfoot
Título & Using SQLite \_ {[}Small\_ Fast\_ Reliable\_ Choose any three \\
Autor & Jay A\_ Kreibich \\
Año & 2010 \\
Editorial & O\textquotesingle Reilly Media, Incorporated \\
Categoría & Lenguajes de Programacion \\
Iniciado & 2026-04-18 \\
Archivo & {[}{[}\url{file:/home/emanuel/Biblioteca/cursos_mit/Electrical} Engineering and Computer Science (Course 6)/02\textsubscript{Programming}\_Languages/Using SQLite \_ {[}Small\_ Fast\_ Reliable\_ Choose any three -- Jay A\_ Kreibich -- 1st ed, Sebastopol, Calif, 2010 -- O\textquotesingle Reilly Media, Incorporated -- isbn13 9780596521189 -- f5659c5ad500c400ff7bd98aa6c304f8 -- Anna's Archive.pdf{]}{[}PDF{]}{]} \\
\end{longtable}
}

\section{Resumen ejecutivo}\label{resumen-ejecutivo}

SQLite es un sistema de gestión de bases de datos relacionales (RDBMS) de dominio público, ligero y autosuficiente, diseñado para almacenar y gestionar datos de manera eficiente sin requerir un servidor separado.

SQLite proporciona una implementación completa de un RDBMS, lo que significa que puede almacenar y gestionar registros definidos por el usuario en grandes tablas. Además, puede procesar comandos de consulta complejos que combinan datos de múltiples tablas para generar informes y resúmenes de datos. La arquitectura de SQLite se basa en los siguientes principios:

\begin{itemize}
\tightlist
\item
  \textbf{\textbf{Serverless}}: SQLite no requiere un proceso de servidor separado o sistema para operar. En su lugar, la biblioteca de SQLite accede directamente a los archivos de almacenamiento. Esto se puede representar como:
\end{itemize}

\[ \text{Aplicación} \leftrightarrow \text{Biblioteca de SQLite} \leftrightarrow \text{Archivo de almacenamiento} \]

\begin{itemize}
\tightlist
\item
  \textbf{\textbf{Zero Configuration}}: Debido a que no hay servidor, no hay configuración necesaria. La creación de una instancia de base de datos SQLite es tan sencilla como abrir un archivo.
\item
  \textbf{\textbf{Cross-Platform}}: La instancia de base de datos completa reside en un solo archivo cross-platform, lo que elimina la necesidad de administración adicional.
\item
  \textbf{\textbf{Autocontenido}}: Una sola biblioteca contiene el sistema de base de datos completo, que se integra directamente en una aplicación host.
\item
  \textbf{\textbf{Pequeña huella de ejecución}}: La construcción predeterminada es menos de un megabyte de código y requiere solo algunos megabytes de memoria. Con algunos ajustes, tanto el tamaño de la biblioteca como el uso de memoria se pueden reducir significativamente.
\end{itemize}

SQLite se encuentra en el contexto de los sistemas de gestión de bases de datos relacionales, que son fundamentales en la informática y las ciencias de la información. Autores clave en este campo incluyen a Edgar F. Codd, quien desarrolló el modelo relacional, y a Larry Ellison, cofundador de Oracle. El debate actual en este campo se centra en la comparación entre los diferentes RDBMS, incluyendo aspectos como el rendimiento, la escalabilidad y la seguridad. SQLite es particularmente popular en aplicaciones embebidas y en dispositivos móviles debido a su pequeño tamaño y bajo consumo de recursos.

\section{Por qué lo leo / Objetivo de lectura}\label{por-quuxe9-lo-leo-objetivo-de-lectura}

\section{Ideas principales}\label{ideas-principales}

\subsection{Idea 1}\label{idea-1}

\subsection{Idea 2}\label{idea-2}

\subsection{Idea 3}\label{idea-3}

\section{Notas por capítulo}\label{notas-por-capuxedtulo}

\subsection{Capítulo 1}\label{capuxedtulo-1}

\subsubsection{Resumen}\label{resumen}

\subsubsection{Conceptos clave}\label{conceptos-clave}

\subsubsection{Preguntas y dudas}\label{preguntas-y-dudas}

\section{Citas y fragmentos destacados}\label{citas-y-fragmentos-destacados}

\section{Vocabulario técnico nuevo}\label{vocabulario-tuxe9cnico-nuevo}

{\def\LTcaptype{none} % do not increment counter
\begin{longtable}[]{@{}lll@{}}
\toprule\noalign{}
Término & Definición & Página \\
\midrule\noalign{}
\endhead
\bottomrule\noalign{}
\endlastfoot
& & \\
\end{longtable}
}

\section{Ejercicios y problemas resueltos}\label{ejercicios-y-problemas-resueltos}

\section{Código relevante}\label{cuxf3digo-relevante}

\begin{Shaded}
\begin{Highlighting}[]
\CommentTok{\# Ejemplo del libro}
\end{Highlighting}
\end{Shaded}

\section{Conexiones con otros libros/conceptos}\label{conexiones-con-otros-librosconceptos}

\section{Aplicaciones prácticas}\label{aplicaciones-pruxe1cticas}

\section{Valoración final}\label{valoraciuxf3n-final}

\begin{description}
\tightlist
\item[Dificultad]
{[} {]} Fácil {[} {]} Medio {[} {]} Difícil {[} {]} Muy difícil
\item[Utilidad]
{[} {]} Baja {[} {]} Media {[} {]} Alta {[} {]} Esencial
\item[Recomendaría]
{[} {]} Sí {[} {]} No {[} {]} Con reservas
\item[Releer]
{[} {]} Sí {[} {]} No
\end{description}

\subsection{Opinión personal}\label{opiniuxf3n-personal}

\section{Anotaciones del PDF}\label{anotaciones-del-pdf}
